% §4 The Index-Set Framework
%
% Presents the SRL framework: update, parity, occupation sets,
% the encoding algorithm, and shows JW and Parity as instances.
% Sets up the tree-derived generalization that the star-tree theorem
% will address.

Seeley, Richard, and Love~\cite{seeley2012} observed that the
Jordan--Wigner encoding can be abstracted: rather than hardcoding the
$Z$-string pattern, one can specify three set-valued functions and
apply a universal formula.  This \emph{index-set framework} is the
most widely cited approach to understanding fermion-to-qubit encodings.

\subsection{The three index sets}
\label{sec:three-sets}

\begin{definition}[Index-set encoding scheme]
\label{def:index-set-scheme}
An \emph{index-set encoding scheme} for $n$ modes is a triple of
set-valued functions $\scheme = (U, P, \mathrm{Occ})$:
\begin{align}
  U &: \{0, \ldots, n{-}1\} \to 2^{\{0,\ldots,n{-}1\}},
    \label{eq:update-fn} \\
  P &: \{0, \ldots, n{-}1\} \to 2^{\{0,\ldots,n{-}1\}},
    \label{eq:parity-fn} \\
  \mathrm{Occ} &: \{0, \ldots, n{-}1\} \to 2^{\{0,\ldots,n{-}1\}}.
    \label{eq:occ-fn}
\end{align}
We call $U(j)$ the \emph{update set}, $P(j)$ the \emph{parity set},
and $\mathrm{Occ}(j)$ the \emph{occupation set} of mode~$j$.
\end{definition}

Informally:
\begin{itemize}
  \item $U(j)$: qubits whose stored parity must be \emph{updated}
    when mode~$j$ changes occupation.
  \item $P(j)$: qubits whose current value encodes the \emph{parity}
    needed to reconstruct mode~$j$'s occupation.
  \item $\mathrm{Occ}(j)$: qubits that collectively store the
    \emph{occupation} of mode~$j$.
\end{itemize}

\subsection{The encoding algorithm}
\label{sec:encoding-algorithm}

Given a scheme~$\scheme$, the Majorana operators for mode~$j$ are:
\begin{align}
  c_j &= X_{\{j\} \cup U(j)} \;\cdot\; Z_{P(j)},
  \label{eq:cj-from-scheme} \\[4pt]
  d_j &= Y_j \;\cdot\; X_{U(j)} \;\cdot\;
    Z_{(P(j)\,\triangle\,\mathrm{Occ}(j))\,\setminus\,\{j\}},
  \label{eq:dj-from-scheme}
\end{align}
where $X_S$ denotes $X$ on every qubit in set~$S$ (identity elsewhere),
$\triangle$ is the symmetric difference, and qubit $j$ receives $X$
in $c_j$ and $Y$ in $d_j$.  The sets $\{j\} \cup U(j)$ and $P(j)$
must be disjoint for the formula to be unambiguous.

The ladder operators then follow from~\eqref{eq:encoded-ladder}.

\begin{remark}[Relation to the SRL formula]
\label{rem:srl-formula}
The original SRL paper~\cite{seeley2012} presents a more complex
formula that places $Y$ at positions in $\mathrm{Occ}(j) \cap P(j)$
and $X$ at $\mathrm{Occ}(j) \setminus P(j)$.  For the three classical
schemes (JW, BK, Parity), both formulas produce identical Majorana
operators.  Our simpler
formulation~\eqref{eq:cj-from-scheme}--\eqref{eq:dj-from-scheme}
suffices for the structural analysis and matches the open-source
implementation~\cite{fockmap2026}.
\end{remark}

\subsection{Jordan--Wigner as an index-set scheme}
\label{sec:jw-as-scheme}

The Jordan--Wigner encoding (\cref{sec:jw}) corresponds to:
\begin{equation}
  U(j) = \emptyset, \quad
  P(j) = \{0, \ldots, j{-}1\}, \quad
  \mathrm{Occ}(j) = \{j\}.
  \label{eq:jw-sets}
\end{equation}

Substituting into~\eqref{eq:cj-from-scheme}:
$c_j = X_j \cdot Z_{\{0,\ldots,j-1\}}$, recovering the JW Majorana
operator.  The update set is empty (each qubit stores exactly one
mode's occupation), and the parity set is the linear chain of all
lower-index qubits.

\subsection{Parity encoding as an index-set scheme}
\label{sec:parity-as-scheme}

The Parity encoding stores cumulative occupation parities:
\begin{equation}
  U(j) = \{j{+}1, \ldots, n{-}1\}, \quad
  P(j) = \{j{-}1\}^*, \quad
  \mathrm{Occ}(j) = \{j{-}1, j\}^*.
  \label{eq:parity-sets}
\end{equation}
(For $j = 0$: $P(0) = \emptyset$, $\mathrm{Occ}(0) = \{0\}$.)

Here the update set propagates changes to all higher-index qubits
(which store cumulative parities), and the parity set is the single
predecessor.  The Parity encoding has the same $O(n)$ weight as
Jordan--Wigner but in the ``opposite direction'': creation operators
act on higher-index qubits rather than lower-index ones.

\subsection{Generalizing to arbitrary trees}
\label{sec:tree-derived}

Both JW and Parity have a natural tree interpretation:
\begin{itemize}
  \item JW corresponds to a \emph{star tree} rooted at a node with
    all modes as children, ordered $0, 1, \ldots, n{-}1$.
  \item Parity corresponds to a star tree with the reverse ordering.
\end{itemize}

This observation suggests a natural generalization: given any labelled
rooted tree $T$ on $[n]$, derive index sets by reading off structural
relationships:
\begin{align}
  U(j) &= \text{ancestors of } j, \notag \\
  P(j) &= \text{remainder}(j) \cup \text{children}(j), \notag \\
  \mathrm{Occ}(j) &= \text{descendants}(j) \cup \{j\},
  \label{eq:tree-derived-sets}
\end{align}
where $\text{remainder}(j)$ collects children of ancestors of $j$
that have index~${} < j$ and are not themselves ancestors of~$j$.

This derivation always produces well-formed sets and yields Pauli
strings of the correct width.  The question is whether the resulting
encoding satisfies the CAR.  We prove in \cref{sec:star-tree} that it
does \emph{if and only if} the tree is a star --- a far sharper
restriction than the framework's generality might suggest.
